\chapter{分物与除法}\label{ch:g3:division}

20 颗弹珠平均分给 4 个孩子，每人分得多少颗？
除法可以回答所有分物问题 --- 也可以回答分组问题 ---
它与乘法互为孪生运算：
二者互相还原。

\section{除法的两种意义}

\begin{definition}[除法]\label{def:g3:division:def}
除法 $20 \div 4$ 可以回答两类问题：
\begin{itemize}
  \item \emph{分物}：$20$ 颗弹珠平均分给 $4$ 个孩子 ---
        \emph{每人}分到多少颗？（$5$ 颗）；
  \item \emph{分组}：$20$ 颗弹珠每袋装 $4$ 颗 ---
        一共能装多少\emph{袋}？（$5$ 袋）。
\end{itemize}
两个问题的答案相同，因为二者还原的是
同一个乘法：$4 \times 5 = 20$。
\end{definition}

\begin{omfigure}
\begin{tikzpicture}[scale=0.55]
  \foreach \g in {0,1,2,3} {
    \draw[omExo, rounded corners] (\g*4.9-0.65,-0.7) rectangle
      (\g*4.9+3.85,0.7);
    \foreach \c in {0,...,4} \fill[omDef] (\g*4.9+\c*0.8,0) circle (0.27);
  }
  \node[below, font=\small] at (8.95,-1.15)
    {$20$ 颗弹珠，分成 $4$ 个相等的组：每组 $5$ 颗};
\end{tikzpicture}

{\small 将 $20$ 平均分成 $4$ 组。若改为每 $4$ 个一组，
则得到 $5$ 组 --- 是同一个除法。}
\end{omfigure}

\begin{method}[用乘法口诀做除法]\label{met:g3:division:tables}
要计算 $42 \div 6$，可以问：“$6$ 乘以\emph{几}得 $42$？”
查 $6$ 的乘法口诀可得：$6 \times 7 = 42$，所以 $42 \div 6 = 7$。
每道除法题，都是把乘法题
倒过来读。
\end{method}

\section{除不尽的情况}

\begin{definition}[余数]\label{def:g3:division:remainder}
把 $23$ 颗弹珠平均分给 $4$ 个孩子，每人分得 $5$ 颗 ---
还剩 $3$ 颗（\emph{剩余}），不够再分一轮。
记作
\[
23 = 4 \times 5 + 3 ,
\]
并称：\emph{商}\index{商}为 $5$，
\emph{余数}\index{余数}为 $3$。余数总是
小于孩子的人数 ---
否则还能再分一轮。
\end{definition}

\begin{omfigure}
\begin{tikzpicture}[scale=0.55]
  \foreach \g in {0,1,2,3} {
    \draw[omExo, rounded corners] (\g*4.9-0.65,-0.7) rectangle
      (\g*4.9+3.85,0.7);
    \foreach \c in {0,...,4} \fill[omDef] (\g*4.9+\c*0.8,0) circle (0.27);
  }
  \foreach \c in {0,1,2} \fill[omThm] (19.7+\c*0.8,0) circle (0.27);
  \node[above, font=\footnotesize, omThm] at (20.5,0.5) {剩余};
  \node[below, font=\small] at (10.2,-1.15)
    {$23 = 4 \times 5 + 3$：商 $5$，余数 $3$（红色）};
\end{tikzpicture}

{\small 把 $23$ 颗弹珠分给 $4$ 人：分 $5$ 轮后，还剩 $3$ 颗
--- 不够再分第 $6$ 轮。}
\end{omfigure}

\begin{example}\label{ex:g3:division:remainder}
计算 $58$ 除以 $7$。查看 $7$ 的乘法口诀：
$7 \times 8 = 56$ 是不超过 $58$ 的最大的积
（$7 \times 9 = 63$ 太大）。所以
\[
58 = 7 \times 8 + 2 :
\]
商 $8$，余数 $2$。验算：$56 + 2 = 58$，且 $2 < 7$。
\end{example}

\begin{method}[解释余数]\label{met:g3:division:interpret}
除完之后，回到题目情境中：
\begin{enumerate}
  \item “每人分得多少？／能装满几组？” --- 商；
  \item “还剩多少？” --- 余数；
  \item “装下全部需要多少个箱子？” --- 商
        若余数不为零，则商\emph{加一}。
\end{enumerate}
\end{method}

\begin{example}\label{ex:g3:division:interpret}
$58$ 名学生去划独木舟，每条独木舟坐 $7$ 人。$58 = 7 \times 8 + 2$：
有八条独木舟坐满，还剩 $2$ 名学生 --- 他们也需要
一条独木舟。因此需要预订 $9$ 条独木舟。
\end{example}

\section{一半、三分之一、四分之一}

\begin{definition}[一半、三分之一、四分之一]\label{def:g3:division:half}
一个数的\emph{一半}\index{一半}是它除以 $2$ 所得的结果；
\emph{三分之一}则除以 $3$；\emph{四分之一}\index{四分之一}则除以 $4$。例如
$18$ 的一半是 $9$，$18$ 的三分之一是 $6$，$36$ 的四分之一
是 $9$。
\end{definition}

\begin{example}[倍与一半相伴]\label{ex:g3:division:double}
求 $46$ 的一半：由于 $46 = 40 + 6$，两部分各取一半，
得 $20 + 3 = 23$。用两倍检验：$23 + 23 = 46$。
熟记两倍数（$25 \to 50$，$45 \to 90$，$150 \to 300$），
求一半便可即刻得出。
\end{example}

\section{练习}

\begin{exercise}[$\star$]\label{exo:g3:division:1}
利用乘法口诀逆向计算：$36 \div 4$；\quad $63 \div 9$；
\quad $40 \div 5$；\quad $56 \div 8$。
\end{exercise}

\begin{exercise}[$\star$]\label{exo:g3:division:2}
对每个除法，写出等式 $a = b \times q + r$，并求出
商和余数：$17 \div 5$；\quad $30 \div 4$；\quad
$50 \div 6$。
\end{exercise}

\begin{exercise}[$\star$]\label{exo:g3:division:3}
$28$ 颗弹珠平均分给 $6$ 个孩子。每人各分到多少颗，
还剩多少颗？如有需要，可画图分一分。
\end{exercise}

\begin{exercise}[$\star$]\label{exo:g3:division:4}
计算：$68$ 的一半；$27$ 的三分之一；$48$ 的四分之一；
$150$ 的一半。
\end{exercise}

\begin{exercise}[$\star$]\label{exo:g3:division:5}
判断下列说法是否正确，并用乘法检验。
“$45 \div 5 = 9$”；\quad “$52 \div 6 = 8$ 余 $4$”；
\quad “$70 \div 8 = 9$ 余 $2$”。
\end{exercise}

\begin{exercise}[$\star$]\label{exo:g3:division:6}
$32$ 张照片贴进相册，每页贴 $4$ 张。需要用多少页，
这是哪一种除法（平均分还是分组）？
\end{exercise}

\begin{exercise}[$\star$]\label{exo:g3:division:7}
一根 $75$ 厘米长的绳子，剪成每段 $9$ 厘米的小段。能剪成多少段完整的，
剩下的一小段有多长？
\end{exercise}

\begin{exercise}[$\star$]\label{exo:g3:division:8}
$50$ 名儿童乘坐每辆有 $4$ 个座位的汽车外出。需要多少辆汽车，
才能坐下？（注意：每个人都要乘车！）
\end{exercise}

\begin{exercise}[$\star$]\label{exo:g3:division:9}
求出下列空缺的数：$? \div 7 = 6$；\quad
$54 \div ? = 6$；\quad $?$ 的一半是 $35$。
\end{exercise}

\begin{exercise}[$\star\star$]\label{exo:g3:division:10}
阿米娜把她的 $47$ 张贴纸尽量公平地分给她的 $5$ 个朋友，
剩下的留给自己。每个朋友各得到多少张贴纸，
阿米娜自己留下多少张？
\end{exercise}

\begin{exercise}[$\star\star$]\label{exo:g3:division:11}
一个数除以 $6$ 时，所有可能的
余数有哪些？除以 $6$ 的商恰为 $7$ 的数中，
最大的数是多少？
\end{exercise}
